%%%%%%%%%%%%%%%%%%%%% chapter-04.tex %%%%%%%%%%%%%%%%%%%%%%%%%%%%%%%%%
%
% Capítulo 4: Isomorfismos y transformaciones topológicas entre espacios celulares
%
%%%%%%%%%%%%%%%%%%%%%%%%%%%%%%%%%%%%%%%%%%%%%%%%%%%%%%%%%%%%%%%%

\chapter{Isomorfismos y transformaciones topológicas entre espacios celulares}
\label{chap:isomorfismos}

\abstract*{En este capítulo se establece la base formal para la equivalencia entre topologías de red, centrando el análisis en la transición de teselaciones hexagonales hacia estructuras ortogonales. Se definen los homomorfismos espaciales necesarios para preservar la dinámica de las reglas de transición, garantizando que las propiedades de computabilidad analizadas en capítulos previos no se vean alteradas por la representación en memoria. A través de una axiomatización geométrica y el desarrollo de ecuaciones de transformación afín, se demuestra la viabilidad del mapeo submatricial de $2 \times 2$ como solución a la discontinuidad de los sistemas de coordenadas discretos.}

\abstract{En este capítulo se establece la base formal para la equivalencia entre topologías de red, centrando el análisis en la transición de teselaciones hexagonales hacia estructuras ortogonales. Se definen los homomorfismos espaciales necesarios para preservar la dinámica de las reglas de transición, garantizando que las propiedades de computabilidad analizadas en capítulos previos no se vean alteradas por la representación en memoria. A través de una axiomatización geométrica y el desarrollo de ecuaciones de transformación afín, se demuestra la viabilidad del mapeo submatricial de $2 \times 2$ como solución a la discontinuidad de los sistemas de coordenadas discretos.}

\section{Axiomatización geométrica de las teselaciones hexagonales y métricas planares}
\label{sec:axiomatizacion_hexagonal}

La discretización del plano euclidiano $\mathbb{R}^2$ mediante polígonos regulares permite la construcción de espacios celulares donde la adyacencia es uniforme. A diferencia de la retícula estándar definida en la Sección 2.1 \eqref{eq:espacio_celular} % TODO: Verificar etiqueta del capítulo 2 sobre definición de espacios celulares
, la teselación hexagonal se caracteriza por poseer el mayor grado de simetría axial en el plano, eliminando las ambigüedades de vecindad presentes en las mallas cuadrangulares (donde se debe distinguir entre adyacencia de arista y de vértice).

\subsection{Representación paramétrica de la celda hexagonal regular}
\label{subsec:parametrizacion}

Se define una celda hexagonal regular $H$ como el conjunto de puntos en el plano acotados por seis segmentos de igual longitud. Sea $s$ la longitud del lado del hexágono, el cual en el modelo unitario se considera $s=1$. Las propiedades geométricas fundamentales que rigen la posición de los vértices y la comunicación entre celdas dependen de dos constantes: el radio circunscrito $R$ y el apotema $a$.

Dado un centro $C = (x_c, y_c)$, los vértices $V_k$ para $k \in \{0, \dots, 5\}$ se determinan mediante las ecuaciones paramétricas:
\begin{equation}
    V_k = \left( x_c + R \cdot \cos\left( \frac{k\pi}{3} \right), y_c + R \cdot \sin\left( \frac{k\pi}{3} \right) \right)
    \label{eq:vertices_hexagono}
\end{equation}

Es posible observar que, por la naturaleza del polígono regular, $R = s$. El apotema, que representa la distancia mínima del centro al punto medio de cada arista, se define como:
\begin{equation}
    a = \frac{\sqrt{3}}{2} s
    \label{eq:apotema}
\end{equation}

Esta configuración asegura que cualquier punto $P \in \mathbb{R}^2$ pueda ser mapeado unívocamente a una celda hexagonal mediante un proceso de cuantización espacial, similar al análisis de complejidad de localización discutido en la Sección \ref{sec:complejidad_renderizado}.

\subsection{Sistemas de coordenadas afines para mallas de grado seis}
\label{subsec:coordenadas_afines}

Para la indexación lógica de las celdas en un autómata celular complejo, el uso de coordenadas cartesianas ortogonales resulta insuficiente debido al desfase intrínseco de las filas. En su lugar, se formaliza el uso de sistemas de coordenadas afines. 

Se define el sistema de coordenadas axiales $(q, r)$, donde los ejes se sitúan a $60^\circ$ de separación. No obstante, para simplificar el cálculo de distancias y algoritmos de colisión de partículas \cite{Basurto2009}, se prefiere la proyección en coordenadas cúbicas $(x, y, z)$ restringidas al plano:
\begin{equation}
    x + y + z = 0
    \label{eq:restriccion_cubica}
\end{equation}

Bajo esta axiomatización, cada celda hexagonal se identifica como un vector en $\mathbb{Z}^3$. Esta representación facilita la definición de invariancia por traslación, propiedad discutida por Kari \cite{Kari2005} como requisito para la universalidad computacional en sistemas dinámicos discretos.

\subsection{Definición formal de vecindad de adyacencia planar $\mathcal{N}_H$}
\label{subsec:vecindad_hexagonal}

Sea $L$ la retícula hexagonal. Se define axiomáticamente la vecindad de adyacencia $\mathcal{N}_H$ de una celda central $c \in L$ como el conjunto de celdas que comparten una frontera unidimensional (arista). A diferencia de la vecindad de Moore en retículas cuadradas \cite{Wolfram2002}, en la topología hexagonal se cumple que:
\begin{equation}
    |\mathcal{N}_H(c)| = 6, \quad \forall c \in L
\end{equation}

La métrica de distancia entre dos celdas $c_1(x_1, y_1, z_1)$ y $c_2(x_2, y_2, z_2)$ en el espacio hexagonal se formaliza como la norma del taxista adaptada (Manhattan hexagonal):
\begin{equation}
    d_H(c_1, c_2) = \frac{|x_1 - x_2| + |y_1 - y_2| + |z_1 - z_2|}{2}
    \label{eq:distancia_hexagonal}
\end{equation}

% TODO: Insertar diagrama que muestre la vecindad hexagonal de radio 1 y la disposición de los ejes axiales y cúbicos.
\begin{figure}[ht]
    \centering
    % \includegraphics[width=0.6\textwidth]{figuras/vecindad_hexagonal.png}
    \caption{Representación de la vecindad $\mathcal{N}_H$ y el sistema de coordenadas cúbicas subyacente.}
    \label{fig:vecindad_hexagonal}
\end{figure}

Esta métrica es fundamental para el análisis de propagación de \textit{gliders} que se abordará en la Parte III, ya que permite predecir la trayectoria de la información sin las distorsiones métricas de las diagonales euclidianas.



\section{Homomorfismos espaciales y condiciones de equivalencia reticular}
\label{sec:homomorfismos_espaciales}

Tras haber definido las propiedades geométricas y la métrica de la teselación hexagonal en la Sección \ref{sec:axiomatizacion_hexagonal}, es imperativo establecer las bases teóricas que justifican la traslación de este continuo discretizado hacia una representación puramente matricial. En el presente apartado se introduce la teoría topológica que legitima el mapeo de un espacio celular a otro, garantizando que la dinámica de la información, la función de transición local $\delta$ y la computabilidad analizadas en el Capítulo \ref{cap:teoria_formal} permanezcan inalteradas bajo la nueva morfología espacial.

\subsection{Isomorfismos sobre grafos reticulares infinitos}
\label{subsec:isomorfismos_grafos}

Para formalizar la equivalencia topológica entre dos espacios celulares distintos, es preciso abstraer la red subyacente empleando la teoría de grafos. Sea $G = (V, E)$ un grafo reticular infinito, donde el conjunto de vértices $V$ representa los índices espaciales de las celdas y el conjunto de aristas $E$ define las relaciones de adyacencia estipuladas por la vecindad axiomática.

Bajo esta premisa, se denota como $G_H = (V_H, E_H)$ al grafo derivado de la teselación hexagonal gobernado por la vecindad $\mathcal{N}_H$, y como $G_Q = (V_Q, E_Q)$ al grafo correspondiente a la retícula ortogonal $\mathbb{Z}^2$ \eqref{eq:espacio_celular} % TODO: Verificar etiqueta del Capítulo 2 para la definición Z^2
. Un homomorfismo espacial entre ambas topologías se establece a través de una función de mapeo $\Phi$ que proyecta los elementos de un dominio sobre el otro preservando sus estructuras de conectividad.

No obstante, debido a la diferencia intrínseca en los grados de los vértices (grado 6 uniforme en $G_H$ frente a grado 4 u 8 en $G_Q$), resulta matemáticamente imposible establecer un isomorfismo biyectivo trivial de cardinalidad uno a uno entre celdas individuales que conserve de manera isótropa las distancias \cite{Toffoli1987}. Por consiguiente, se define un mapeo biyectivo hacia subredes, donde cada vértice $v \in V_H$ se asocia unívocamente a un subgrafo inducido $B \subset V_Q$. Para que el sistema celular se considere topológicamente equivalente y capaz de simular la misma física discreta, la función $\Phi$ debe cumplir la condición de adyacencia: si $(u, v) \in E_H$, entonces debe existir una adyacencia efectiva y computable entre los bloques de subconjuntos $\Phi(u)$ y $\Phi(v)$ en el dominio ortogonal.

\subsection{Teorema de recubrimiento y descomposición celular ortogonal}
\label{subsec:teorema_recubrimiento}

El núcleo analítico de la transformación geométrica propuesta en este documento reside en el axioma de equivalencia espacial mediante bloques matriciales de cardinalidad superior. Se establece la partición cuadrangular extendida como la solución formal a la discontinuidad geométrica descrita anteriormente.

Se formaliza el axioma de que el área, la capacidad de almacenamiento de estado y, crucialmente, las direcciones de adyacencia de una celda hexagonal $h \in V_H$, pueden ser representadas de forma isomórfica mediante un bloque ortogonal contiguo de celdas cuadradas. Específicamente, el mapeo $\Phi(h)$ proyecta la celda sobre una subred $B_{2\times 2}$ definida como:
\begin{equation}
    B_{2\times 2} = \left\{ c_{i,j}, c_{i+1,j}, c_{i,j+1}, c_{i+1,j+1} \right\} \subset V_Q
    \label{eq:bloque_2x2}
\end{equation}

Esta relación biyectiva, donde un único hexágono encapsula la información espacial de cuatro cuadrados subyacentes, responde a una necesidad estrictamente topológica. Como es posible observar, una celda ortogonal individual posee únicamente 4 aristas disponibles (vecindad de Von Neumann), lo cual es topológicamente deficiente para enrutar de manera simétrica las 6 trayectorias de propagación de partículas abordadas en la dinámica de colisiones \cite{Basurto2009}. 

Al expandir la unidad fundamental de cómputo a un bloque de $2 \times 2$, el perímetro externo de esta nueva meta-celda expone exactamente 8 segmentos ortogonales. Este incremento en los grados de libertad perimetrales garantiza el isomorfismo geométrico: 6 de las caras del bloque cuadrangular se asignan biyectivamente a las 6 vecindades hexagonales evaluadas por la métrica \eqref{eq:distancia_hexagonal}, asegurando que ninguna trayectoria se superponga prematuramente.

% TODO: Insertar diagrama conceptual del mapeo. Debe mostrar un hexágono abstracto superpuesto transformándose en un bloque de 2x2 cuadrados, ilustrando cómo las 6 aristas del hexágono se asientan sobre el perímetro de los 4 cuadrados.
\begin{figure}[ht]
    \centering
    % \includegraphics[width=0.7\textwidth]{figuras/isomorfismo_2x2.png}
    \caption{Equivalencia topológica: Mapeo isomórfico suprayectivo entre una celda hexagonal con grado de adyacencia 6 y una partición cuadrangular de $2 \times 2$.}
    \label{fig:isomorfismo_2x2}
\end{figure}

Esta descomposición celular asegura además que el modelo computacional derivado en la Sección 3.5 posea un tiempo de indexación en memoria estrictamente $\mathcal{O}(1)$, eludiendo las aberraciones de redondeo propias de las coordenadas de punto flotante y conservando el rigor exigido por la teoría clásica de autómatas celulares \cite{Kari2005}.